% Part: many-valued-logic
% Chapter: three-valued-logics
% Section: kleene

\documentclass[../../../include/open-logic-section]{subfiles}

\begin{document}

\olfileid{mvl}{thr}{skl}

\olsection{ক্লিনি যুক্তিবিদ্যাসমূহ}

স্টিফেন ক্লিনি দুটি ত্রিমানী যুক্তিবিদ্যা প্রবর্তন করেন। এগুলির প্রেরণা ছিল
এমন এক যুক্তিবিদ্যা, যেখানে সত্যমানগুলিকে গণনামূলক প্রক্রিয়ার ফল হিসেবে
ভাবা হয়: কোনো প্রক্রিয়া~$\True$ বা~$\False$ দিতে পারে, আবার থামতে
ব্যর্থও হতে পারে। শেষ ক্ষেত্রে সংশ্লিষ্ট সত্যমানটি অসংজ্ঞায়িত; একে
সত্যমান~$\Undef$ দিয়ে প্রকাশ করা হয়।

কোনো বচন~$!A$-এর নঞর্থক রূপ গণনা করতে প্রথমে~$!A$-এর মান গণনা করবে,
তারপর ফলটির বিপরীত মান ফেরত দেবে। $!A$-এর গণনা না থামলে গোটা
প্রক্রিয়াটিও থামবে না; তাই $\Undef$-এর নঞর্থক মান~$\Undef$।

কোনো সংযোজন $!A \land !B$ গণনা করার দুটি উপায় আছে। প্রথমে~$!A$,
তারপর~$!B$ গণনা করা যায়; উভয়ের ফল~$\True$ হলে ফল হবে~$\True$, অন্যথায়
$\False$। কোনো একটি গণনা থামতে ব্যর্থ হলে গোটা প্রক্রিয়াটিও থামবে না।
তাই এই ক্ষেত্রে একটি সংযোজ্য অসংজ্ঞায়িত হলে সংযোজনটিও অসংজ্ঞায়িত।
বিযোজনের ক্ষেত্রেও একই কথা প্রযোজ্য।

কিন্তু $!A$ ও~$!B$-কে সমান্তরালে মূল্যায়ন করতে পারলে আমরা আরও ভালো ফল
পেতে পারি। তখন প্রক্রিয়া দুটির একটি থেমে~$\False$ ফেরত দিলেই আমরা থামতে
পারি, কারণ উত্তরটি মিথ্যা হতেই হবে। এই ক্ষেত্রে তাই একটি সংযোজ্য মিথ্যা
হলে অন্যটি অসংজ্ঞায়িত হলেও সংযোজনটি মিথ্যা। একইভাবে, কোনো বিযোজন
সমান্তরালে গণনা করার সময় বিযোজ্য দুটির একটির প্রক্রিয়া সত্য ফেরত দিলেই
আমরা থামতে পারি; তখন বিযোজনটি সত্য হতেই হবে। ফলে কোনো উপাদান
অসংজ্ঞায়িত হলেও যৌগিক উক্তির ফল কী হবে তা এই ক্ষেত্রে জানা যায়। এই
ব্যাখ্যায়~$\Undef$-কে “অসংজ্ঞায়িত”-এর বদলে “অজ্ঞাত” বলে পড়া যেতে পারে।

এই দুই ব্যাখ্যা ক্লিনির শক্তিশালী ও দুর্বল যুক্তিবিদ্যার জন্ম দেয়।
শর্তবচনকে $\lnot !A \lor !B$-এর সমতুল্য হিসেবে সংজ্ঞায়িত করা হয়।

\begin{defn}
\emph{শক্তিশালী ক্লিনি যুক্তিবিদ্যা}~$\LogKs$ নিচের ম্যাট্রিক্স দিয়ে
সংজ্ঞায়িত:
\begin{enumerate}
  \item $\lnot$, $\land$, $\lor$, $\lif$-সহ প্রমিত বচনগত ভাষা
  $\Lang L_0$।
  \item সত্যমানের সেট $V = \{\True, \Undef, \False\}$।
  \item $\True$ একমাত্র মনোনীত মান; অর্থাৎ $V^+ = \{\True\}$।
  \item সত্যমান-অপেক্ষকগুলি নিচের সারণিগুলি দিয়ে দেওয়া হয়:
  \begin{center}
    \begin{tabular}{c|c}
      $\tf{\lnot}$ & \\
      \hline
      $\True$ & $\False$ \\
      $\Undef$ & $\Undef$ \\
      $\False$ & $\True$
    \end{tabular}
    \quad
    \begin{tabular}{c|ccc}
      $\tf{\land}[\LogKs]$ & $\True$ & $\Undef$ & $\False$ \\
      \hline
      $\True$ & $\True$ & $\Undef$ & $\False$ \\
      $\Undef$ & $\Undef$ & $\Undef$ & $\False$\\
      $\False$ & $\False$ & $\False$ & $\False$
    \end{tabular}
    \\[2ex]
    \begin{tabular}{c|ccc}
      $\tf{\lor}[\LogKs]$ & $\True$ & $\Undef$ & $\False$ \\
      \hline
      $\True$ & $\True$ & $\True$ & $\True$ \\
      $\Undef$ & $\True$ & $\Undef$ & $\Undef$ \\
      $\False$ & $\True$ & $\Undef$ & $\False$
    \end{tabular}
    \quad
    \begin{tabular}{c|ccc}
      $\tf{\lif}[\LogKs]$ & $\True$ & $\Undef$ & $\False$ \\
      \hline
      $\True$ & $\True$ & $\Undef$ & $\False$ \\
      $\Undef$ & $\True$ & $\Undef$ & $\Undef$  \\
      $\False$ & $\True$ & $\True$ & $\True$
    \end{tabular}
  \end{center}
\end{enumerate}
\end{defn}

\begin{defn}
\emph{দুর্বল ক্লিনি যুক্তিবিদ্যা}~$\LogKw$ নিচের ম্যাট্রিক্স দিয়ে
সংজ্ঞায়িত:
\begin{enumerate}
  \item $\lnot$, $\land$, $\lor$, $\lif$-সহ প্রমিত বচনগত ভাষা
  $\Lang L_0$।
  \item সত্যমানের সেট $V = \{\True, \Undef, \False\}$।
  \item $\True$ একমাত্র মনোনীত মান; অর্থাৎ $V^+ = \{\True\}$।
  \item সত্যমান-অপেক্ষকগুলি নিচের সারণিগুলি দিয়ে দেওয়া হয়:
  \begin{center}
    \begin{tabular}{c|c}
      $\tf{\lnot}$ & \\
      \hline
      $\True$ & $\False$ \\
      $\Undef$ & $\Undef$ \\
      $\False$ & $\True$
    \end{tabular}
    \quad
    \begin{tabular}{c|ccc}
      $\tf{\land}[\LogKw]$ & $\True$ & $\Undef$ & $\False$ \\
      \hline
      $\True$ & $\True$ & $\Undef$ & $\False$ \\
      $\Undef$ & $\Undef$ & $\Undef$ & $\Undef$\\
      $\False$ & $\False$ & $\Undef$ & $\False$
    \end{tabular}
    \\[2ex]
    \begin{tabular}{c|ccc}
      $\tf{\lor}[\LogKw]$ & $\True$ & $\Undef$ & $\False$ \\
      \hline
      $\True$ & $\True$ & $\Undef$ & $\True$ \\
      $\Undef$ & $\Undef$ & $\Undef$ & $\Undef$ \\
      $\False$ & $\True$ & $\Undef$ & $\False$
    \end{tabular}
    \quad
    \begin{tabular}{c|ccc}
      $\tf{\lif}[\LogKw]$ & $\True$ & $\Undef$ & $\False$ \\
      \hline
      $\True$ & $\True$ & $\Undef$ & $\False$ \\
      $\Undef$ & $\Undef$ & $\Undef$ & $\Undef$  \\
      $\False$ & $\True$ & $\Undef$ & $\True$
    \end{tabular}
  \end{center}
\end{enumerate}
\end{defn}

\begin{prop}
  $\LogKs$ ও~$\LogKw$-এ কোনো সর্বতঃসত্য নেই।
\end{prop}

\begin{proof}
  প্রতিটি !!{propositional variable}s~$p$-এর জন্য
  $\pAssign v(p) = \Undef$ হলে যেকোনো সূত্র~$!A$-এর সত্যমান হবে
  $\pValue v(!A) = \Undef$; কারণ উভয় যুক্তিবিদ্যায়
  \[
    \tf{\lnot}(\Undef) = \tf{\lor}(\Undef, \Undef) = \tf{\land}(\Undef,
  \Undef) = \tf{\lif}(\Undef, \Undef) = \Undef
  \]
  $\LogKs$ ও~$\LogKw$—কোনোটির জন্যই $\Undef \notin V^+$; তাই এই
  !!{valuation}-এ $!A$ মনোনীত হবে না।
\end{proof}

দুর্বল ও শক্তিশালী ক্লিনি যুক্তিবিদ্যা—উভয়েরই কোনো সর্বতঃসত্য না থাকলেও
তাদের অ-তুচ্ছ অনুসিদ্ধান্ত সম্পর্ক আছে।

\begin{prob}
  নিচের কোন সম্পর্কগুলি (ক)~শক্তিশালী এবং (খ)~দুর্বল ক্লিনি যুক্তিবিদ্যায়
  সিদ্ধ? প্রতিটির জন্য একটি সত্যসারণি দাও।
  \begin{enumerate}
    \item $p, p \lif q \Entails q$
    \item $p \lor q, \lnot p \Entails q$
    \item $p \land q \Entails p$
    \item $p \Entails p \land p$
    \item $p \Entails p \lor q$
  \end{enumerate}
\end{prob}

দিমিত্রি বোচভার $\Undef$-কে “অর্থহীন” হিসেবে ব্যাখ্যা করেন এবং
কূটাভাসী বাক্যগুলিকে~$\Undef$ মান দেওয়ার বিধান করে মিথ্যাবাদীর
কূটাভাসের মতো কূটাভাস সমাধানে এটি ব্যবহারের চেষ্টা করেন। তিনি এমন একটি
যুক্তিবিদ্যা প্রবর্তন করেন, যা মূলত অতিরিক্ত কয়েকটি সংযোজক দিয়ে প্রসারিত
দুর্বল ক্লিনি যুক্তিবিদ্যা। এগুলির দুটি হলো “বহিঃস্থ নঞর্থকরণ” ও
“অসংজ্ঞায়িত” অপারেটর:
\begin{center}
  \begin{tabular}{c|c}
    $\tf{\sim}$ & \\
    \hline
    $\True$ & $\False$ \\
    $\Undef$ & $\True$ \\
    $\False$ & $\True$
  \end{tabular}
\quad
  \begin{tabular}{c|c}
    $\tf{+}$ & \\
    \hline
    $\True$ & $\False$ \\
    $\Undef$ & $\True$ \\
    $\False$ & $\False$
  \end{tabular}
\end{center}

\begin{prob}
  বোচভারের যুক্তিবিদ্যায় $\sim$-কে কি $\lnot$ ও~$+$ দিয়ে সংজ্ঞায়িত
  করতে পারো? অর্থাৎ শুধু !!{propositional
  variable}~$p$ আছে এবং
  $\sim$ নেই—এমন !!a{formula} খুঁজে বের করো, যার সত্যমান সবসময়
  $\mathord{\sim}p$-এর সত্যমানের সমান। তোমার উত্তর সঠিক দেখাতে একটি
  সত্যসারণি দাও।
\end{prob}


\end{document}
